\section{A15.1b2a: mixed strip $R \ge e$ injectivity}
\label{sec:p12-b2a}

\begin{theorem}[Mixed strip, high source radius]\label{thm:p12-b2a}
Let $2a < T_0 < c$, $e \le R < T$, $T < S < T_0$.  Then
\(\ker L_{R,S,T_0}^{\{a,b,2a\}} = \{0\}\).
\end{theorem}

\begin{proof}
Let $h \in L^2(R, S)$ with $L_{T_0} h = 0$.  Set $H(t) := h(T + t)$ for
$0 < t < \sigma$, $\sigma := S - T$.

With the newly opened source horizon, the E-source $u = a + x$
introduces exactly one new forward term, giving
\[
p\, h(x) + r\, \mathrm{sgn}(x - d)\, h(|x - d|) - q\, h(a - x) - p\, H(x)\,\mathbf 1_{x < \sigma}
\;=\; 0,
\qquad 0 < x < a.
\tag{E$_\sigma$}
\]

We show $h = 0$ on the lower half $(0, a)$ first, then $H = 0$.

\medskip
\noindent\emph{Step 1 (upper interior: $d < x < a$).}
Here $0 < x - d < a - d = e \le R$ and $0 < a - x < a - d = e \le R$, so
both auxiliary values are outside the support.  (E$_\sigma$) reduces to
$p\, h(x) = 0$ (the $H(x)\mathbf 1_{x<\sigma}$-term also vanishes because
$x > d > \sigma$).  Since $p > 0$,
\[
h(x) = 0 \quad\text{a.e. on }(d, a).
\tag{S1}
\]

\medskip
\noindent\emph{Step 2 (lower interior: $R < x < d$).}
Here $0 < d - x < d - R \le d - e = \delta < e \le R$, so
$h(d - x) = 0$.  Equation (E$_\sigma$) reduces to
\[
p\, h(x) - q\, h(a - x) = 0
\qquad\text{whenever the $H$-term vanishes.}
\tag{2}
\]
The $H$-term vanishes for $x \ge \sigma$; and $\sigma < \varepsilon < e
\le R < x$, so $x > \sigma$ throughout $(R, d)$.

If $a - x \le R$, immediately $h(x) = 0$.

If $a - x > R$: since $R \ge e$, we have $a - x < a - R \le a - e = d$, so
$a - x \in (R, d)$.  Applying (2) at $a - x$ gives
\[
p\, h(a - x) - q\, h(x) = 0.
\tag{3}
\]
Combining (2) and (3),
\[
(p^2 - q^2)\, h(x) = 0.
\]
Since $(p/q)^2 = 2^{3/2} \ne 1$, we have $p^2 \ne q^2$, so
\[
h(x) = 0 \quad\text{a.e. on }(R, d).
\tag{S2}
\]

Together with (S1) and the boundary continuity, this gives
\(h = 0\) a.e. on $(0, a)$.

\medskip
\noindent\emph{Step 3 (upper half via A14.3a upper source).}
For $u = 3a - t$ (or the equivalent upper E-source, in the notation
inherited from A14.3a §2.2), the upper-half fold structure and the
same $(p^2 - q^2)$-mechanism kill $h$ on $(a, T)$.  This uses only
$R \ge e$ combined with (S2) and requires no new terms from the
$2a$-shift beyond those already accounted for.  We omit the parallel
argument; it is a direct transcription of A14.3a §5--6 with $S \le T$
replaced by $S \le T_0 < c < 3a$.

Consequently $h = 0$ on $(0, T)$.

\medskip
\noindent\emph{Step 4 (tail dies).}
For $0 < t < \sigma$, (E$_\sigma$) at $x = t$ gives
\[
p\, h(t) + r\, \mathrm{sgn}(t - d)\, h(|t - d|) - q\, h(a - t) - p\, H(t) = 0.
\]
Since $t < \sigma < \varepsilon < e < d$, we have $t < d$ and $\mathrm{sgn}(t-d) = -1$.
The three arguments $t$, $d - t$, $a - t$ all lie in $(0, a)$ and are
therefore in the null region proved in Step~3.  The equation reduces to
\[
-p\, H(t) = 0,
\]
whence $H(t) = 0$ on $(0, \sigma)$, i.e. $h = 0$ on $(T, S)$.
\end{proof}

\begin{remark}[Numerical cross-check]\label{rem:p12-b2a-numerical}
A finite-dimensional Galerkin discretization of $L_{R,S,T_0}$ with
$R \in [0.144, 0.30]$, $S = 0.5$, $T_0 = 0.78$ and $N = 200$ mesh
points yields mesh-stable positive $\sigma_{\min}(L_{R,S,T_0})$,
consistent with Theorem~\ref{thm:p12-b2a}.
\end{remark}
